\def\xlong{20}
\def\ylong{10}
\def\Tmax{0.9*\ylong}
\def\ti{0.35*\xlong}
\def\tf{0.65*\xlong}

\begin{tikzpicture}[
    texto/.style={black, font=\sffamily},
    lim ancho/.style={text width=2.5cm, align=center},
    eje/.style={thick, arrows=-Stealth},
    limites/.style={dashed, DarkRed},
    marca/.style={loosely dotted, DarkRed},
    linea central/.style={ultra thick, DarkBlue}
]

\draw[eje] (0,0) -- (\xlong,0) node[midway, below=5pt, texto] {Tiempo (min)};
\draw[eje] (0,0) -- (0,\ylong) node[pos=0.6, rotate=90, left=10pt, texto] {Temperatura (\unit{\kelvin})};

\draw[linea central] (0,1) to[out=60,in=200] (\ti,\Tmax) -- (\tf,\Tmax) to[out=-60,in=160] (\xlong,1);

\draw[limites] (\ti,0) -- (\ti,1.1*\Tmax);
\draw[limites] (\tf,0) -- (\tf,1.1*\Tmax);

\draw[marca] (0,\Tmax) -- (\ti,\Tmax) node[midway, above, texto, lim ancho] {Temperatura de sostenimiento};

\node[texto, lim ancho] at (0.6*\ti, 0.4*\Tmax) {Fase de calentamiento};
\node[texto, lim ancho] at (0.5*\xlong, 0.6*\Tmax) {Fase de sostenimiento};
\node[texto, lim ancho] at (0.5*\xlong+0.6*\tf, 0.6*\Tmax) {Fase de enfriamiento};

\end{tikzpicture}